\section{MCMC 在无法独立采样时构造平稳链}
\label{sec:13-Machine-Learning/08-Sampling-and-Inference/02}

你给一个药物剂量反应模型做了 Bayesian 拟合，六个参数、三层结构，似然和先验都写在代码里，每个 \(\theta\) 的后验密度值随手就能算出来。可你要的是后验均值和 \(95\%\) 可信区间，而这需要从后验里采样——偏偏没有任何现成的随机数生成器认识这个分布。上一节的 Monte Carlo 假设你能从 \(p\) 独立抽样，这个假设在这里直接落空。

障碍出在归一化常数上。后验

\[
\pi(\theta)=p(\theta\given y)
=\frac{p(y\given\theta)p(\theta)}{p(y)},
\qquad
p(y)=\int p(y\given\theta)p(\theta)\,d\theta,
\]

分子逐点可算，分母是一个六维积分，跟上一节那个电网积分一样算不动。于是你知道 \(\pi\) 的形状，却不知道它的高度。逆变换采样要分布函数，拒绝采样要一个能罩住 \(\pi\) 的上界，两者都用不上。

MCMC 的做法是不去求 \(\pi\)，而是造一条以 \(\pi\) 为平稳分布的 Markov 链，让它自己走出来。链的定义、转移核与多步演化在 \hyperref[sec:06-Random-Processes/03-Markov-Chains/01]{``Markov 性质、转移矩阵与多步演化''} 里已经建好，这里只是把状态空间换成参数空间，把要求换成一条：链的长期访问频率必须是 \(\pi\)。

\subsection{细致平衡是设计原则，不是待验证的性质}

如果按常规思路，要造一条平稳分布是 \(\pi\) 的链，得先写下转移核再解方程 \(\pi K=\pi\)。这是个不动点问题，在连续状态空间上并不比原积分容易。Metropolis--Hastings 反过来做：先随手挑一个提议分布，再用一个接受概率把它掰成满足细致平衡的样子。

链的状态记为 \(x\)（数据 \(y\) 已固定，不再是变量）。当前在 \(x\)，从提议分布 \(q(x'\given x)\) 抽出候选 \(x'\)，以

\[
\alpha(x,x')
=\min\set{1,\ \frac{\pi(x')\,q(x\given x')}{\pi(x)\,q(x'\given x)}}
\]

接受，否则原地不动。记 \(K(x,x')\) 为一步转移核（含被拒时停在原处的点质量），直接代入可验证

\[
\pi(x)K(x,x')=\pi(x')K(x',x),
\]

两边对 \(x\) 积分，右端 \(K\) 的积分为一，立刻得到 \(\pi K=\pi\)。骨架就这么短：细致平衡是逐对状态的等式，平稳性是它的积分推论，方向由前者到后者，反过来不成立——细致平衡是充分条件而非必要条件，不满足它的链也可以有正确的平稳分布。

接受率里那个比值才是整件事能落地的关键。\(\pi\) 的未知归一化常数 \(p(y)\) 在分子分母上原样约掉，你只需要能算未归一化的 \(p(y\given\theta)p(\theta)\)。写不出的那个积分从此不必写。所有需要归一化常数的量——后验的绝对密度值、模型证据——仍然算不出，MCMC 绕过的是采样，不是配分函数本身，这一点留到本单元最后一节再算账。

\subsection{合法性来自遍历性，与样本独立无关}

平稳性只说明：若当前状态已经服从 \(\pi\)，下一步仍服从 \(\pi\)。它没说从你随手挑的初值出发能不能到达 \(\pi\)。补上这一步的是遍历性——不可约、非周期、常返，这些条件在 \hyperref[sec:06-Random-Processes/03-Markov-Chains/04]{``平稳分布、可逆性与遍历收敛''} 里有完整讨论。条件满足时，遍历定理给出

\[
\frac1n\sum_{t=1}^nh(X_t)\longrightarrow\E_\pi[h(X)]
\]

几乎处处成立。注意这条定理里没有出现``独立''二字。时间平均之所以收敛到空间平均，是因为链把整个状态空间按 \(\pi\) 的比例访问了一遍，而不是因为相邻样本互不相关。相邻样本当然相关，而且往往高度相关——这不影响估计的相合性，只影响它的精度。

精度的账要单独算。记 \(\sigma_h^2=\Var_\pi(h(X))\)，\(\rho_k=\Corr(h(X_t),h(X_{t+k}))\)，则

\[
\Var(\bar h_n)
\approx\frac{\sigma_h^2}{n}\Bigl(1+2\sum_{k\ge1}\rho_k\Bigr)
=\frac{\sigma_h^2\tau}{n},
\]

括号里的 \(\tau\) 是积分自相关时间，链相当于只提供了 \(n/\tau\) 个独立样本。\(\tau=200\) 时，一百万次迭代折合五千个独立样本，这就是有效样本量的含义。

收敛快慢由转移算子的谱决定。把 \(K\) 看成作用在函数上的线性算子，它的最大特征值是一，对应常函数；次大特征值的模 \(\lambda_2\) 与一之差称为谱间隙，初始分布偏离 \(\pi\) 的部分大致按 \(\lambda_2^t\) 衰减，同时 \(\tau\approx(1+\lambda_2)/(1-\lambda_2)\)。特征值支配演化速率这件事在 \hyperref[sec:03-Linear-Algebra/05-Eigenvalues-and-Eigenvectors/02]{``对角化与动力系统：把耦合演化变成独立缩放''} 里已经讲清楚，这里只是同一个机制换了个场合。由此可见 burn-in 长度与自相关时间不是两件事：谱间隙小则初始瞬态消退慢、相邻样本相关强，两个症状同源。指望靠加长 burn-in 修好一条自相关时间为几万的链，是在同一堵墙上换个地方撞。

提议尺度的调节也从这里得到解释。随机游走步长过小，接受率接近一，链却几乎原地踏步，\(\lambda_2\) 贴近一；步长过大，候选频频落进低密度区被拒，链同样不动。接受率只是可观察的线索，效率本身要看目标函数的有效样本量与 Monte Carlo 标准误。

\subsection{链没跨过模态时你拿到的是一条被困住的轨道}

多峰后验把上面所有讨论推到极限。设想两个后验峰之间隔着一片极低密度的谷地，随机游走要跨过去需要连续若干次都往同一个方向走且不被拒绝，概率小到实践中等于零。此时链在理论上仍然不可约（谷地密度大于零），遍历定理仍然成立，只是收敛所需的时间尺度远超你的算力预算。

\begin{center}
\begin{tikzpicture}[>=stealth]
\begin{scope}
\draw[->] (-3.9,0) -- (4.0,0) node[right]{\(x\)};
\draw[thick] plot[smooth] coordinates
  {(-3.6,0.03) (-3.0,0.14) (-2.6,0.50) (-2.2,1.05) (-2.0,1.25)
   (-1.8,1.05) (-1.4,0.50) (-1.0,0.14) (-0.6,0.05) (0,0.02)
   (0.6,0.05) (1.0,0.14) (1.4,0.50) (1.8,1.05) (2.0,1.25)
   (2.2,1.05) (2.6,0.50) (3.0,0.14) (3.6,0.03)};
\foreach \p in {-2.42,-2.18,-1.94,-2.30,-2.06,-1.78,-2.54,-1.90,-2.26,-2.10}
  \fill (\p,-0.42) circle (1.3pt);
\node at (-2.15,-0.95) {链访问过的区域};
\node[gray] at (2.05,-0.95) {从未到达};
\node at (0,1.65) {\(\pi(x)\)};
\end{scope}
\begin{scope}[shift={(5.9,0)}]
\draw[->] (0,-0.45) -- (5.4,-0.45) node[right]{\(t\)};
\draw[->] (0,-0.55) -- (0,1.75) node[above]{\(x_t\)};
\draw[gray,dashed] (0,1.25) -- (5.2,1.25);
\node[gray] at (4.1,1.5) {另一个模态};
\draw[thick] plot[smooth] coordinates
  {(0.15,0.52) (0.5,0.28) (0.9,0.44) (1.3,0.22) (1.7,0.46)
   (2.1,0.30) (2.5,0.48) (2.9,0.26) (3.3,0.42) (3.7,0.31)
   (4.1,0.47) (4.5,0.25) (4.9,0.40) (5.3,0.33)};
\node at (2.7,-0.95) {单链 trace 平稳得毫无破绽};
\end{scope}
\end{tikzpicture}
\end{center}

\emph{右图的平稳只说明链没离开左峰；它没有发现右峰，也就无从表现出不安。}

图里那条 trace 是最容易骗过人的证据。它没有趋势、没有跳变、方差稳定，任何单链诊断都会给它开绿灯，而它对应的估计把后验质量的一半判为不存在。删掉前一半样本毫无帮助——burn-in 删掉的是同样困在左峰里的点，它不会凭空造出一次跨峰移动。

能识破这种失败的诊断只有一类：跨链比较。从彼此分散的初值并行跑若干条链，比较链间方差与链内方差，这正是 \(\widehat R\) 做的事。所有链困在同一个峰时 \(\widehat R\) 也接近一，所以初值必须真正分散；一旦有链落到不同峰上，\(\widehat R\) 会明显大于一并报警。配合每个参数的有效样本量、自相关函数和 Monte Carlo 标准误一起看，才构成一份能用的诊断表。

诊断的逻辑方向要摆正：它们是反证工具。链间不一致是明确的失败证据；指标全部正常只说明在这次运行的算力范围内没发现问题，不能证明后验已被走遍。这个不对称是有限时间运行的固有属性，任何统计量都改不了。

改善探索的手段往往在参数化里。高度相关的参数让随机游走沿着一条狭长的脊来回蹭，层级模型的漏斗形几何在颈部把步长要求压得极小；中心化与非中心化之间换一种写法、把各参数标准化到同一尺度，常常比调采样器参数有效得多。利用梯度的提议——Hamiltonian Monte Carlo 及其自适应版本——能沿着等高线走出很长的一步，把 \(\tau\) 降一两个数量级，但它同样跨不过高势垒，多峰问题依旧要靠回火、并行温度或多起点来处理。

最后区分两件容易混起来的事。采样诊断说的是``链有没有正确地表示这个后验''，后验预测检验说的是``这个模型能不能生成像真实数据那样的数据''。前者全绿而后者失败，说明模型错了而采样对了，此时加长链没有任何用处。下一节要走的是另一条完全不同的路：不再让样本去逼近后验，而是换一族简单分布把后验的位置直接优化出来。
